\section{Синтаксический анализатор PTX}

\subsection{Требования к парсеру и выбор стратегии}

\subsubsection{Анализ структуры PTX-текста}

PTX-текст имеет строго линейную структуру: последовательность директив,
за которыми следуют тела функций. Каждое тело функции~--- последовательность
базовых блоков; каждый базовый блок~--- последовательность инструкций,
предваряемых необязательной меткой. В отличие от языков программирования
общего назначения, PTX не содержит вложенных выражений или скобочных
структур внутри отдельных инструкций.

Эти свойства обосновывают отказ от построения полного AST и использование
однопроходного линейного разбора. Сложность: $O(N)$ по числу строк $N$,
что критически важно для PTX-файлов размером 5--50~КБ, возникающих
при встраивании трансцендентных функций компилятором NVCC.

\subsubsection{Однопроходный инкрементный разбор}

Парсер реализует конечный автомат с тремя состояниями:
$Q = \{q_{\mathrm{init}}, q_{\mathrm{func}}, q_{\mathrm{block}}\}$
и следующими переходами:
\begin{itemize}
  \item $q_{\mathrm{init}} \xrightarrow{\texttt{.entry/.func}} q_{\mathrm{func}}$:
        встретили заголовок функции.
  \item $q_{\mathrm{func}} \xrightarrow{\$\mathrm{label}:} q_{\mathrm{block}}$:
        встретили метку базового блока.
  \item $q_{\mathrm{block}} \xrightarrow{\mathrm{instruction}} q_{\mathrm{block}}$:
        добавляем инструкцию в текущий блок.
  \item $q_{\mathrm{block}} \xrightarrow{\}} q_{\mathrm{init}}$:
        закрывающая скобка тела функции.
\end{itemize}

На каждом шаге поддерживается контекст: текущая функция, текущий
базовый блок, таблица файлов \texttt{source\_file}/\texttt{source\_line},
монотонный счётчик PC.

Таким образом, однопроходный инкрементный разбор обеспечивает
линейную сложность $O(N)$ по числу строк PTX-текста и полное
отсутствие возвратов. Результатом является структурированный объект
\texttt{PtxModule}, содержащий таблицу инструкций с монотонными
PC-адресами, граф базовых блоков и таблицу символов для разрешения
переходов во время исполнения.

\subsection{Реализация парсера}

\subsubsection{Разбор объявлений функций и параметров}

Заголовок функции разбирается регулярным выражением:

\begin{lstlisting}[language=C++, caption={Регулярное выражение для заголовка функции}]
static const std::regex kFuncRe{
    R"(\.(?:entry|func)\s+)"          // kind
    R"((?:\(.*?\)\s*)?)"              // optional return type
    R"((\w+))"                         // name
    R"(\s*\(([^)]*)\))",              // params
    std::regex::optimize
};
\end{lstlisting}

Результат~--- объект \texttt{Function} с таблицей параметров
\texttt{params\_} и флагом \texttt{isEntry()}, по которому диспетчер
различает точки входа ядра и вспомогательные функции. Параметры
функции разбираются рекурсивно: каждый элемент вида
\texttt{.param .align~N .bM name} порождает запись
\texttt{ParamSlot\{offset, size\}}, используемую исполнителем
при копировании аргументов вызова в регистровый файл.

Директива \texttt{.reg .type \%prefix<N>} объявляет $N$ регистров
заданного типа. Парсер сохраняет:
\texttt{reg\_decls[func][prefix] = (count, type)}.
Эти данные используются исполнителем при инициализации регистрового
файла варпа~--- выделяется ровно столько слотов, сколько объявлено
в PTX, без избыточного выделения. Профилировщик дополнительно
использует эту таблицу для подсчёта \texttt{write\_ops} по префиксу.

\subsubsection{Таблица меток и граф базовых блоков}

Таблица \texttt{BasicBlockAddrMap} является ключевой структурой,
связывающей парсер с исполнителем управляющего потока. Метка вида
\texttt{\$BB0\_3:} фиксирует точку входа в базовый блок:

\begin{enumerate}
  \item Закрываем текущий базовый блок, сохраняем его в \texttt{PtxModule}.
  \item Записываем \texttt{basic\_blocks\_["\$BB0\_3"] = current\_pc}~---
        отображение имени метки в монотонный PC первой инструкции блока.
  \item Открываем новый базовый блок.
\end{enumerate}

При исполнении инструкции ветвления \texttt{bra \$L} исполнитель
вызывает \\
\texttt{GetBasicBlockOffset("\$L")} и устанавливает PC варпа
в полученное значение. Без этой таблицы реализация переходов
потребовала бы линейного поиска по всей таблице инструкций.

\subsubsection{Разбор отдельной инструкции}

Каждая строка-инструкция разбирается в пять последовательных этапов,
полностью определяющих данные, необходимые исполнителю во время
выполнения:

\begin{enumerate}
  \item \textbf{Предикатный префикс.} Опциональный \texttt{@\%pN}
        или \texttt{@!\%pN} извлекается первым и сохраняется
        в поле \texttt{predicate}. Исполнитель проверяет его перед
        выполнением: если предикатный регистр ложен, инструкция
        пропускается без изменения состояния регистрового файла.
  \item \textbf{Нормализация мнемоники.} \texttt{fma.rn.f32}
        разбивается на ключ поиска в таблице диспетчеризации (\texttt{fma})
        и список модификаторов \texttt{["rn", "f32"]}, которые передаются
        исполнителю для выбора режима округления и ширины типа.
  \item \textbf{Список операндов.} Строка разбивается по запятой
        с учётом скобок \texttt{[...]}, порождая упорядоченный вектор
        \texttt{Operand}. Порядок критичен: первый операнд~--- регистр
        назначения для большинства инструкций.
  \item \textbf{Классификация и разрешение операндов}
        (подробно~--- в следующем разделе).
  \item \textbf{Присвоение монотонного PC.} Счётчик инкрементируется
        для каждой исполняемой инструкции и используется как ключ
        в \texttt{instruction\_table\_} и в записях \texttt{profiling.txt}.
\end{enumerate}

\subsubsection{Разрешение операндов}

Классификация каждого операнда производится единожды при разборе,
а не при каждом исполнении инструкции, что устраняет накладные
расходы на разбор строк в горячем цикле исполнителя:

\begin{itemize}
  \item \textbf{Именованный регистр}: \texttt{\%r5}, \texttt{\%f12},
        \texttt{\%rd0}, \texttt{\%p3}. Префикс однозначно определяет тип
        (\texttt{r}~$\to$~\texttt{s32}, \texttt{f}~$\to$~\texttt{f32},
        \texttt{rd}~$\to$~\texttt{u64}, \texttt{p}~$\to$~\texttt{pred})
        и индекс слота в регистровом файле варпа.
  \item \textbf{Числовой литерал}: десятичный (\texttt{42}),
        шестнадцатеричный (\texttt{0x3f800000}) или IEEE~754
        (\texttt{0f3f800000}). Значение конвертируется в \texttt{uint64\_t}
        при разборе и хранится непосредственно в \texttt{Operand}.
  \item \textbf{Адресное выражение}: \texttt{[\%rd5+8]}~--- базовый
        регистр и смещение разделяются при разборе; исполнитель
        вычисляет адрес как \texttt{reg[base] + imm} без дополнительного
        разбора строки.
  \item \textbf{Символическая метка}: \texttt{\$BB0\_3}~---
        хранится как строка; при первом исполнении перехода
        разрешается через \texttt{BasicBlockAddrMap} и кэшируется.
\end{itemize}

\subsection{Внутреннее представление инструкций}

\subsubsection{Иерархия классов инструкций и монотонный счётчик PC}

Каждая PTX-инструкция представлена отдельным классом, наследующим
абстрактный базовый класс:

\begin{lstlisting}[language=C++, caption={Базовый класс (build/app/ptx\_asm/autogen/instructions.h)}]
class Instruction {
public:
    virtual void Execute(std::shared_ptr<WarpContext>& wc) = 0;
    virtual std::string_view Name() const = 0;
    virtual bool HasExplicitPredicate() const { return false; }
};
\end{lstlisting}

Конкретные подклассы генерируются из \texttt{instructions.json}.
Каждый хранит только поля, специфичные для своей мнемоники
(пространство памяти, тип данных, режим округления, операнды):

\begin{lstlisting}[language=C++, caption={Пример: класс ldInstruction}]
class ldInstruction : public Instruction {
public:
    ldspaceQl   space_;   // Global, Shared, Local, Const
    dataType    data_;    // F32, F64, U32, ...
    addressOperand addr_; // base register + immediate offset
    registerOperand dst_; // destination register
    std::optional<registerOperand> pred_; // predicate

    void Execute(std::shared_ptr<WarpContext>& wc) override;
    std::string_view Name() const override { return "ld"; }
    bool HasExplicitPredicate() const override;
};
\end{lstlisting}

Монотонный PC присваивается инструкции при разборе PTX и хранится
в \texttt{InstructionList} модуля как индекс в глобальном векторе:
\texttt{Module::GetInstruction(pc)} возвращает \texttt{shared\_ptr<Instruction>}
за $O(1)$. Тот же PC используется как ключ в записях \texttt{profiling.txt},
обеспечивая инъективное соответствие между трассой и таблицей инструкций.

При исполнении \texttt{Module::GetInstruction(pc)} возвращает
\texttt{const Instruction\&} за O(1) через \texttt{instruction\_table\_[pc]}.

\newpage
